\chapter*{LAMPIRAN}
%\addcontentsline{toc}{chapter}{LAMPIRAN}	

\section{Isi Lampiran}

Lampiran bersifat opsional bergantung hasil kesepakatan dengan pembimbing 
dapat berupa:

\begin{enumerate}
\item Bukti pelaksanaan Kuesioner seperti pertanyaan kuesioner, resume jawaban 
responden, dan dokumentasi kuesioner. 
\item Spesifikasi Aplikasi atau Sistem yang dikembangkan meliputi spesifikasi 
teknis aplikasi, tautan unduh aplikasi, manual penggunaan aplikasi, hingga 
screenshot aplikasi. 
\item Cuplikan kode yang sekiranya penting dan ditambahkan. 
\item Tabel yang terlalu panjang yang masih diperlukan tetapi tidak 
memungkinkan untuk ditayangkan di bagian utama skripsi.
\item Gambar-gambar pendukung yang tidak terlalu penting untuk ditampilkan di 
bagian utama. Akan tetapi, mendukung argumentasi/pengamatan/analisis.
\item Penurunan rumus-rumus atau pembuktian suatu teorema yang terlalu
panjang dan terlalu teknis sehingga Anda berasumsi bahwa pembaca biasa
tidak akan menelaah lebih lanjut. Hal ini digunakan untuk memberikan
kesempatan bagi pembaca tingkat lanjut untuk melihat proses penurunan
rumus-rumus ini.
\end{enumerate}

\section{Rumus Algoritma Tata Letak Graf \textit{Force-Directed}}
\label{lamp:rumus-graf}
Bagian ini memuat persamaan matematis dari ketiga gaya yang digunakan pada algoritma tata letak graf \textit{force-directed} pada visualisasi jaringan sanad, beserta nilai parameternya.

\begin{enumerate}
\item \textbf{Gaya tolak (\textit{repulsion})}, bekerja antar pasangan node dan hanya diterapkan jika jarak antar node kurang dari 320 piksel:
$$ F_{tolak} = \frac{k_r}{d^2}, \quad k_r = 2000 $$
dengan $d$ adalah jarak Euclidean antara dua node.

\item \textbf{Gaya pegas (\textit{spring})}, bekerja pada setiap edge dengan panjang ideal 200 piksel:
$$ F_{pegas} = (d - L_0)\cdot k_s, \quad L_0 = 200,\ k_s = 0{,}003 $$

\item \textbf{Gaya radial (\textit{radial constraint})}, menarik setiap node menuju jari-jari cincin targetnya ($r_{target}$) sesuai peringkat sentralitas:
$$ F_{radial} = (r - r_{target})\cdot k_c, \quad k_c = 0{,}014 $$
dengan $r$ adalah jarak node saat ini dari titik pusat.
\end{enumerate}

Setelah ketiga gaya dijumlahkan, kecepatan tiap node diperbarui lalu dikalikan dengan faktor peredam (\textit{damping}) sebesar 0,78. Posisi awal node ditentukan berdasarkan peringkat sentralitasnya pada cincin konsentris dengan jari-jari 180 (peringkat 1--5), 320 (peringkat 6--15), 460 (peringkat 16--30), dan 600 (sisanya).